% ============================================================================
% PLANTILLA 11 — BANCO DE PREGUNTAS
% ----------------------------------------------------------------------------
% Para: bancos de preguntas organizados por tema y nivel cognitivo
% (Taxonomía de Bloom), con claves visibles para el docente.
% Compila por defecto CON claves; para la versión sin claves cambie la
% primera línea por \documentclass{evaluacion}.
% Compilar: latexmk -pdf plantilla-11-banco-de-preguntas.tex
% ============================================================================
\documentclass[claves]{evaluacion}

\universidad{Universidad Nacional de San Cristóbal de Huamanga}
\facultad{Facultad de Ciencias Económicas, Administrativas y Contables}
\escuela{Escuela Profesional de Economía}
\curso{Macroeconomía I}
\codigocurso{EC-243}
\docente{Mg. Nombre Apellido}
\periodo{Actualizado a julio de 2026}
\tipoevaluacion{Banco de Preguntas}
\fechaevaluacion{2026-07-21}
\modalidad{Documento interno del docente}
\autorevaluacion{Edison Achalma}

\begin{document}
\portada[Unidades 1--2 · Contabilidad nacional y modelo IS-LM\\
         Versión con claves — uso exclusivo del docente]

\begin{instrucciones}[Organización del banco]
Cada pregunta indica entre corchetes el nivel cognitivo de la Taxonomía
de Bloom: \textbf{[R]}ecordar, \textbf{[C]}omprender, \textbf{[A]}plicar,
\textbf{[An]}alizar, \textbf{[E]}valuar, \textbf{[Cr]}ear. La numeración
es corrida para poder citar preguntas por código
(p.~ej.\ \enquote{B-\thepregunta}) al armar un examen.
\end{instrucciones}

% ======================================================================
\section{Contabilidad nacional}

\begin{pregunta}[][R · Opción múltiple]
¿Cuál de los siguientes componentes \textbf{no} forma parte del PBI por
el método del gasto?
\begin{opciones}
  \opcion Consumo privado.
  \opcion Inversión bruta interna.
  \opcion \correcta{Compra de un bono del tesoro por un hogar.}
  \opcion Exportaciones netas.
\end{opciones}
\clave{c — es una transacción financiera, no producción corriente}
\end{pregunta}

\begin{pregunta}[][C · Verdadero/Falso]
Indique verdadero o falso:
\begin{tablavf}
  \vf[V]{El PBI nominal puede crecer aunque el PBI real caiga.}
  \vf[F]{El deflactor del PBI y el IPC miden exactamente la misma
         canasta.}
  \vf[V]{La depreciación explica la diferencia entre producto bruto y
         producto neto.}
\end{tablavf}
\end{pregunta}

\begin{pregunta}[][A · Cálculo]
Una economía produce 100 quintales de papa a S/~80 y 40 fanegas de maíz a
S/~120 en el año base; en el año corriente produce 110 y 38 a precios de
S/~95 y S/~130. Calcule el PBI nominal, el PBI real y el deflactor del
año corriente.
\begin{solucion}
Nominal: $110(95)+38(130) = 10\,450 + 4\,940 = 15\,390$.
Real (precios base): $110(80)+38(120) = 8\,800+4\,560 = 13\,360$.
Deflactor: $15\,390/13\,360 = 1{,}1519$ (115,2).
\end{solucion}
\end{pregunta}

\begin{pregunta}[][An · Desarrollo]
La economía informal supera el 70\,\% del empleo en el Perú. Analice dos
consecuencias de este hecho para la medición del PBI y una estrategia que
emplea el INEI para mitigarlas.
\begin{solucion}
Puntos esperados: subestimación del nivel del PBI (producción no
registrada) y sesgo en la composición sectorial; volatilidad espuria
cuando la formalización desplaza actividad ya existente hacia registros.
Mitigación: métodos indirectos (encuestas de hogares ENAHO, coeficientes
insumo-producto, consumo eléctrico) para imputar la producción informal.
\end{solucion}
\end{pregunta}

% ======================================================================
\section{El modelo IS-LM}

\begin{pregunta}[][R · Completar]
La curva IS recoge los equilibrios del mercado de
\completar[bienes]{2.8cm} y la curva LM los del mercado de
\completar[dinero]{2.8cm}; la pendiente de la IS es
\completar[negativa]{2.8cm}.
\end{pregunta}

\begin{pregunta}[][A · Problema]
Considere $C = 40 + 0{,}8(Y - T)$, $I = 150 - 10r$, $G = T = 100$,
demanda real de dinero $L = 0{,}2Y - 8r$ y oferta real $M/P = 96$.
\begin{apartados}
  \item Derive las ecuaciones IS y LM. \pts{2}
  \item Halle el equilibrio $(Y^{*}, r^{*})$. \pts{2}
  \item Calcule el nuevo equilibrio si $G$ aumenta a 120 y cuantifique el
        efecto expulsión. \pts{2}
\end{apartados}
\begin{solucion}
\textbf{a)} IS: $Y = C+I+G = 40+0{,}8Y-80+150-10r+100
\Rightarrow Y = 1\,050 - 50r$.
LM: $0{,}2Y - 8r = 96 \Rightarrow Y = 480 + 40r$.

\textbf{b)} $1\,050 - 50r = 480 + 40r \Rightarrow r^{*} = \frac{570}{90}
= 6{,}33$; $Y^{*} = 480+40(6{,}33) = 733{,}3$.

\textbf{c)} Con $G=120$: IS: $Y = 1\,150 - 50r$; igualando,
$r^{*} = 7{,}44$, $Y^{*} = 777{,}8$. El multiplicador con tasa fija
daría $\Delta Y = 5\times20 = 100$; el aumento efectivo es $44{,}4$:
el efecto expulsión absorbe $55{,}6$ (la subida de $r$ en 1,11 puntos
reduce la inversión en 11,1).
\end{solucion}
\end{pregunta}

\begin{pregunta}[][E · Ensayo corto]
Evalúe la utilidad del modelo IS-LM para analizar la política monetaria
peruana actual, considerando que el BCRP fija la tasa de interés y no la
oferta monetaria. ¿Qué modificación del modelo responde a esta crítica?
\begin{solucion}
Se espera la crítica de endogeneidad del dinero: con meta de tasa, la LM
se vuelve horizontal al nivel de la tasa de política (modelo IS-MP de
Romer). El estudiante debe reconocer que las conclusiones cualitativas de
la política fiscal se mantienen, pero la transmisión monetaria se lee
directamente como elección de $r$ y regla de reacción (tipo Taylor).
\end{solucion}
\end{pregunta}

\begin{pregunta}[][Cr · Diseño]
Diseñe un ejercicio numérico original de estática comparativa en IS-LM
donde una política fiscal expansiva no altere el producto. Especifique
los parámetros que lo hacen posible y la intuición económica.
\begin{solucion}
Respuesta modelo: demanda de dinero insensible a la tasa
($L = kY$, LM vertical): todo aumento de $G$ eleva $r$ hasta expulsar
exactamente la misma cantidad de inversión (efecto expulsión total).
Se acepta cualquier parametrización coherente; se valora la intuición
monetarista del caso.
\end{solucion}
\end{pregunta}

\findelexamen
\end{document}
